\documentclass[10pt,a4paper]{amsart}
\usepackage[margin=19mm]{geometry}
\usepackage{amsmath,amssymb,mathtools}
\usepackage[scr=rsfs,cal=euler]{mathalfa}
\usepackage{xfrac}
\usepackage[unicode,hidelinks,bookmarks=false]{hyperref}
\newcommand{\ie}{i.e.\ }
\newcommand{\cf}{cf.\ }
\newcommand{\resp}{respectively\ }
\newcommand{\Subsection}{\subsection}
\providecommand{\nobreakdash}{\nobreak\hbox{-}}
\setlength{\emergencystretch}{2em}
\multlinegap0pt
\usepackage{amssymb}
\newtheorem{lem}{Lemma}
\title{A simple separable C*-algebra which is not singly generated}
\author{George A. Elliott, Chun Guang Li, Zhuang Niu}
\date{Source: arXiv:2608.13707v1 (CC BY 4.0)}
\begin{document}
\maketitle

\noindent\emph{Source and licence.} A simple separable C*-algebra which is not singly generated. Version arXiv:2608.13707v1 (CC BY 4.0), distributed by arXiv under the Creative Commons Attribution 4.0 International licence, http://creativecommons.org/licenses/by/4.0/, as recorded in the arXiv licence metadata for this exact version and shown on the abstract page https://arxiv.org/abs/2608.13707v1. Full licence notice: \texttt{licenses/arxiv-2608.13707-CC-BY-4.0.txt}.\par\smallskip

\noindent\textit{Editorial scope.} Counted unit: the article's lem nenzero-dc-lem (source0.tex lines 460--462). The article's complete original proof of it is delivered with it (source0.tex lines 464--485, one \texttt{proof} environment of 22 lines). No mathematics is written, repaired or re-derived: the statement and the proof are the article's own lines, in the article's own notation, and no retained line is substituted or re-flowed.  No further source-local context is needed: every cross-reference of every retained line resolves inside the two delivered spans.  Nothing was omitted from any retained span, so the package carries no omission record.  No retained line contains a citation, so the wrapper carries no bibliography and no bibliography entry.  No retained line contains a figure environment, an image-inclusion command or drawing code, so no image is delivered and the package declares no asset dependency.  Editorial formatting only, in addition: an AMS \texttt{amsart} preamble that restates the article's own declarations and macros used by the retained lines, namely the environments \texttt{lem} on separate counters, together with the standard packages those lines need.  The retained environments print in this document as Lemma 1 in order of appearance, whereas the article prints the same items with the numbering of its own full document.  The complete native source of the article as distributed in the arXiv e-print for this exact version is bundled unchanged as \texttt{sources/207605/source0.tex}.
\begin{lem}\label{nenzero-dc-lem}
Let $M$ and $N$ be differentiable manifolds with the same dimension, and assume that $N$ is connected. Let $L \subseteq N$ be a connected differentiable submanifold, and let $f: M \to N$ be a differentiable map which transverses to $L$. Assume that $\mathrm{deg}(\pi) = 1$. Then $\mathrm{deg}(f |_{f^{-1}(L)}) \neq 0$.
\end{lem}

\begin{proof}
Since $f$ transverses to $L$, the preimage $f^{-1}(L)$ is a submanifold of $M$ with same dimension as $L$. Fix an orientation of $f^{-1}(L)$, and then fix an orientation of $M$ which is compatible with $f^{-1}(L)$. Also fix an orientations of $L$ and $N$ which are compatible.

Pick a regular value $y \in L \subseteq N$ for $f: f^{-1}(L) \to L$. Then $f^{-1}(y)$ is a zero-dimensional submanifold of $M$, and hence $$f^{-1}(y) = \{x_1, x_2, ..., x_n\}.$$ 

Consider the tangent spaces $$\mathrm{T}_M(x_i), \quad i=1, ..., n, $$ and the tangent maps
$$\mathrm{d}f(x_i): \mathrm{T}_M(x_i) \to \mathrm{T}_N(y). $$

Then
$$
\mathrm{deg}(f) = \sum_{i=1}^n w_i,
$$
where $w_i = \pm 1$, $i=1, ..., n$, depending on whether the tangent map $\mathrm df: \mathrm{T}_M(x_i) \to \mathrm{T}_N(y)$ preserves the orientation or reverses it. Since $\mathrm{deg}(f) = 1$, necessarily $n$ is odd.

On the other hand
$$
\mathrm{deg}(f|_{f^{-1}(L)}) = \sum_{i=1}^n w'_i,
$$
where $w'_i = \pm 1$, $i=1, ..., n$, depends on whether the tangent map $\mathrm d(f|_{f^{-1}(L)}): \mathrm{T}_{f^{-1}(L)}(x_i) \to \mathrm{T}_L(y)$ preserves the orientation or reverses the orientation. Since $n$ is odd, the sum $\sum_{i=1}^n w'_i$ must be nonzero, as desired.
%
% \mathrm{T}_M(x_2), ..., \mathrm{T}_M(x_n). $$
\end{proof}

\end{document}
